\chapter{\MakeUppercase{Discussão}}\label{sec:discussao}

Este capítulo interpreta os resultados do \autoref{sec:resultados} em conjunto, extrai implicações para a pesquisa em inteligência artificial aplicada a mercados agrícolas, para os produtores e para quem constrói ferramentas de apoio à decisão, e discute as escolhas de projeto que os resultados sugerem revisar.

\section{As Questões de Pesquisa em Conjunto}\label{sec2:disc-rqs}

Lidos isoladamente, os resultados da RQ1 seriam modestos para um trabalho de previsão: a rede perde para o passeio aleatório no horizonte de um mês, perde para a curva de futuros em três meses na validação, e sua vantagem só é inequívoca no horizonte de seis meses. Lidos em conjunto com a RQ2, eles contam outra história. Uma previsão que acerta a direção em quatro de cada cinco casos no horizonte de seis meses, ainda que erre a magnitude, é suficiente para que um motor de decisão sequencial gere ganho financeiro sobre as heurísticas do produtor: 7,8\,\% sobre a venda imediata no teste cego e 6,1\,\% em quatro safras, com significância estatística nas duas amostras. A pergunta que o trabalho responde não é se a rede prevê bem, mas se a informação que ela extrai, por fraca que seja, tem valor econômico quando acoplada a uma regra de decisão que respeita custos e incerteza. A resposta, nas quatro safras analisadas, é afirmativa, com as ressalvas de amostra discutidas adiante.

Três mecanismos explicam como a estrutura de decisão sequencial mitiga os erros da previsão.

O primeiro é a \textbf{assimetria entre acerto e erro}. Quando o motor prevê recuperação a partir de um fundo e acerta, o ganho é de 10 a 30 \% em poucos meses, como em julho e agosto de 2025 e no segundo semestre de 2023. Quando prevê alta e erra, a política rolante reavalia a cada mês e vende assim que a previsão do mês seguinte deixa de justificar o carregamento, o que limita a perda na maioria dos casos. A exceção é setembro de 2022, com perda de 26 \%, quando a rede manteve a previsão de alta por dois meses consecutivos durante uma queda persistente; nesse caso a política rolante não protege, e a perda é da mesma ordem do ganho nos acertos. A assimetria média é, portanto, favorável, mas não garantida. O caso da safra 2024/25 mostra o mecanismo pelo lado positivo: em nenhuma das datas de julho a outubro de 2024 a previsão de seis meses apontava a magnitude da alta, mas a previsão de um mês foi positiva em cada uma delas, e a reavaliação mensal manteve o estoque até que ela deixasse de ser, capturando 34 a 40 \% em cada data. O motor não previu a alta; ele a seguiu mês a mês.

O segundo é a \textbf{incerteza explícita}. O motor não recebe apenas o ponto central da previsão, mas também a dispersão histórica dos resíduos por horizonte. Quando a deriva prevista é pequena diante dessa dispersão, o valor esperado de armazenar mal supera o custo de carregamento, e qualquer aversão a risco leva à venda. É por isso que, com $\lambda = 0{,}25$, a frequência de armazenar cai de 40 para 30 \% das datas sem que a receita média caia: as decisões descartadas eram as de menor convicção e menor ganho. O mesmo mecanismo, porém, torna a política frágil a partir de $\lambda = 0{,}5$: como a dispersão residual do modelo final é pequena, a penalização de risco passa a superar a deriva esperada em quase todas as datas e o motor converge para a venda imediata. O sistema converte a qualidade da previsão em largura do cone de preços, e o cone em prudência, mas o parâmetro de aversão a risco precisa ser calibrado com o produtor em uma faixa estreita, abaixo de 0,5 nesta configuração.

O terceiro é a \textbf{estrutura de gatilhos}. A saída do motor não é uma data de venda, mas uma regra: em cada mês futuro, um preço acima do qual vender é ótimo. Isso torna a política robusta a erros de magnitude da previsão, porque a venda é acionada pela realização do preço e não pela previsão dele. Um produtor que segue a regra vende quando o mercado confirma a expectativa, e não quando a rede a formula.

A RQ3 fecha o argumento. A política é insensível aos três parâmetros econômicos que a literatura de comercialização agrícola recomenda ao produtor acompanhar de perto: Selic, tarifa de armazenagem e base. Isso não significa que esses custos não importem, mas que, no regime de preços do período, a deriva prevista quando o motor decide armazenar é uma ordem de grandeza maior que o custo mensal de carregamento. Em um regime de preços baixos e juros altos a conclusão pode mudar, e o desenho do motor permite testá-lo sem alteração de código. O único parâmetro que altera a política é a aversão a risco, que é justamente o parâmetro que pertence ao produtor e não ao mercado.

Há, por fim, uma razão estrutural para que armazenar na colheita tenda a pagar, e ela ajuda a explicar por que a política ganha mesmo com previsões modestas. O preço do arábica tem sazonalidade: na mediana dos anos civis completos de 2015 a 2025, junho e julho ficam 3\,\% abaixo da média do ano, agosto 1\,\% abaixo, e novembro e dezembro 3 e 4\,\% acima, com os demais meses próximos da média. O vale coincide com a colheita, quando a oferta é abundante, e o pico com o fim do ano, quando ela se esgota. A média mensal amplifica esse padrão, com dezembro 11\,\% acima, porque a alta de 2024 dobrou o preço dentro de um mesmo ano civil, o que mostra também o quanto um único ano pode distorcer a leitura. Uma saca armazenada da colheita até o verão captura, em um ano típico, esse diferencial de 6 a 7 pontos percentuais, que supera com folga o custo de carregamento de cerca de 1\,\% ao mês. A rede aprende esse padrão, e a bienalidade e a codificação do mês são as duas variáveis de maior importância por permutação; o motor, por sua vez, só o explora quando a previsão de curto prazo não contradiz a sazonal, o que é exatamente o que se espera de uma política que combina conhecimento estrutural com informação corrente.

\section{Onde a Ferramenta Falha e Por Quê}\label{sec2:disc-falhas}

A honestidade exige tratar os resultados negativos com o mesmo cuidado dos positivos, e o trabalho tem três.

O primeiro é a simulação por safra. Quando o produtor é modelado como alguém que colhe uma vez por ano e decide a partir dali, a venda parcelada supera o AgroVisible nas duas safras cegas, 1\,837 contra 1\,807 R\$/sc, e a razão está em uma única decisão: em 1º de julho de 2024, com previsão de 21 dias negativa, o motor vendeu a saca inteira a 1\,362 R\$/sc, dois meses antes de uma alta de 84 \%. No \textit{backtest} mensal, a mesma data é apenas uma de vinte, e as dezenove seguintes compensam o erro; na simulação, a saca vendida não volta. O mecanismo que faz a política funcionar, a reavaliação mensal de um estoque existente, só opera enquanto há estoque, e uma venda total na colheita o desliga. A implicação de projeto é direta: a venda total não deveria ser a única alternativa à armazenagem total. Uma política que vende uma fração e reavalia o restante, isto é, que incorpora a venda parcelada ao motor em vez de tratá-la só como referência, preservaria a opção de participar de altas não previstas. Ela não foi implementada neste trabalho.

O segundo é a compressão das previsões. Treinada com erro quadrático em uma série próxima de um passeio aleatório, a rede aprende que a previsão de menor risco é próxima de zero, e só se afasta dela quando o sinal é forte. Em um movimento sem precedente na amostra de treino, como a alta de 2024, o sinal de seis meses não existia: as previsões de 126 dias com origem em julho a outubro de 2024 eram próximas de zero ou negativas, contra retornos realizados acima de 50 \%. O motor capturou a alta pela previsão de um mês e pela reavaliação, não porque a rede a tenha antecipado. Isso sugere prever o excedente sobre a curva de futuros em vez do retorno absoluto, o que faria a rede aprender apenas o que o mercado ainda não precifica e removeria da tarefa dela a parte que a curva já resolve.

O terceiro é a queda persistente. Em setembro e outubro de 2022, e de novo em junho de 2023, a rede previu alta no meio de uma queda que durou meses, e a política rolante, que protege contra erros isolados, não protege contra um erro repetido. A perda de setembro de 2022, de 26 \%, é da mesma ordem dos maiores ganhos. Uma revisão possível é ponderar os horizontes na interpolação pela acurácia direcional de cada um, dando ao horizonte de 126 dias, o mais confiável, maior peso na deriva dos primeiros meses; outra é condicionar a decisão de armazenar à concordância entre a previsão de curto prazo e a de seis meses. Nenhuma das duas foi implementada neste trabalho, e todas as revisões deste item são candidatas naturais a trabalhos futuros.

\section{Implicações para a Pesquisa}\label{sec2:disc-pesquisa}

O resultado central do trabalho é uma evidência empírica de que métricas estatísticas de erro são insuficientes para avaliar sistemas de inteligência artificial voltados a decisões em mercados agrícolas. Três observações sustentam essa afirmação.

Primeiro, o \textit{ranking} dos modelos por RMSE e o \textit{ranking} por valor econômico não coincidem. O ARIMA e o passeio aleatório têm o menor RMSE em 21 dias na validação, mas não carregam informação alguma para a decisão. A LSTM tem RMSE pior nesse horizonte e ainda assim gera ganho, porque acerta a direção em 63 e 126 dias. Um estudo que se encerrasse na tabela de métricas de um mês concluiria pelo modelo errado.

Segundo, a métrica que melhor antecipa o valor econômico não é o erro médio, mas a acurácia direcional no horizonte relevante para a decisão. A ablação mostra isso com clareza: os grupos de variáveis cuja remoção mais eleva o erro de 126 dias, calendário, mercado e retornos, são também os que derrubam a direção nesse horizonte, de 81 para 62 a 72 \%, enquanto a remoção de qualquer grupo é indiferente em 21 dias. A análise exploratória conta a mesma história antes do treino: a correlação das variáveis de oferta, base e calendário com o alvo cresce com o horizonte, e as de curto prazo são ruído em todos eles.

Terceiro, a avaliação econômica exige um protocolo próprio, com contabilidade de custos, referências que o agente decisor já possui e horizonte de teste que cubra mais de um regime de mercado. Os quatro trabalhos do núcleo sobre café mapeados no \autoref{sec:fundamentacao} não atendem a nenhum desses três requisitos, e a lacuna que este trabalho aponta é metodológica antes de ser algorítmica. A proposta é que estudos de previsão de preços agrícolas reportem, ao lado do RMSE, o resultado de pelo menos uma regra de decisão simples aplicada às previsões, contra pelo menos uma heurística de mercado e contra a curva de futuros quando ela existir.

Há também uma implicação sobre baselines. A curva de futuros ajustada pela base sazonal é uma referência exigente que nenhum dos trabalhos mapeados utiliza, e que na validação supera a rede neural no horizonte de três meses. Para commodities com mercado a termo líquido, essa deve ser a referência mínima de qualquer modelo preditivo, porque é a expectativa que o produtor já tem de graça.

\section{Implicações para os Produtores}\label{sec2:disc-produtores}

Para o cafeicultor, o AgroVisible oferece quatro coisas que as heurísticas usuais não oferecem.

A primeira é \textbf{disciplina}. A venda parcelada é uma boa heurística em altas e uma má heurística em quedas, e o produtor não sabe de antemão em qual regime está. A política do motor armazena quando há razão para esperar recuperação ou continuidade de alta e vende quando não há, e no \textit{backtest} de quatro safras foi a melhor estratégia não oracular em três delas e no total, perdendo apenas na safra de queda contínua, em que vender de imediato era a única política correta. A ressalva está na simulação por safra: se o produtor vende tudo na colheita por recomendação do motor, a disciplina vira rigidez, e por isso a recomendação de venda total em julho é a que mais merece ser confrontada com a leitura do próprio produtor.

A segunda é \textbf{proteção contra estocagem com margem negativa}. Ao descontar armazenagem, quebra e custo de oportunidade em cada aresta do grafo, o motor só recomenda carregar o estoque quando a deriva esperada supera o custo total de carregamento. O produtor que armazena por hábito ou por expectativa vaga não faz essa conta.

A terceira é \textbf{leitura da margem sobre o custo de produção}. Com o COE da CONAB por safra, a ferramenta informa em cada decisão a margem sobre o custo operacional efetivo e classifica o resultado esperado nas faixas de lucro supernormal a falência da gestão de custos agrícola \cite{AgroSafety2026}. No período do teste, a margem mínima foi de 74 \% e a restrição de piso não atuou; em 2023, com preços entre 800 e 1\,200 R\$/sc, ela teria atuado em mais de um terço das datas. É nos períodos de preço baixo que essa leitura importa, e é exatamente nesses períodos que o produtor tende a vender por necessidade de caixa.

A quarta é \textbf{redução da assimetria de informação}. O produtor negocia com cooperativas e tradings que acompanham a curva de futuros, a posição dos fundos e o câmbio diariamente. O AgroVisible coloca essas mesmas informações, sintetizadas em uma regra de gatilho por mês, à disposição de quem vende uma saca.

Duas advertências acompanham essas implicações. A ferramenta não substitui a leitura do produtor sobre sua própria situação de caixa, e a restrição de piso sobre o COE foi mantida desligada justamente porque forçar o carregamento quando o preço está abaixo do custo pode ser pior que vender. E a política é calibrada para o mercado físico de referência, não para o preço de balcão de uma cooperativa específica, cuja base local pode diferir da modelada.

\section{Implicações para Construtores de Ferramentas}\label{sec2:disc-tools}

A experiência de construir o AgroVisible sugere sete diretrizes de engenharia para plataformas de apoio à decisão voltadas a cooperativas e empresas agrícolas.

\begin{alineas}
	\item \textbf{Desacoplar previsão e decisão.} A interface entre os dois módulos é um vetor de médias e um de desvios por horizonte. Isso permitiu trocar o alvo da rede, retreiná-la com sementes diferentes e rodar ablações sem tocar no motor, e permitirá substituir a LSTM por qualquer estimador futuro. Ferramentas que embutem a previsão na regra de decisão perdem essa flexibilidade.
	\item \textbf{Exigir incerteza do módulo preditivo.} Um ponto central sem dispersão obriga o motor a tratar todas as previsões como igualmente confiáveis. A dispersão residual por horizonte é barata de calcular e é o que torna a aversão a risco operacional.
	\item \textbf{Entregar regras, não datas.} O gatilho de venda por mês é mais útil e mais robusto que uma recomendação pontual, porque o produtor o executa contra o preço observado. A interface da ferramenta deve exibir a regra, o cone de preços e a distribuição do mês de venda, e não apenas a ação de hoje.
	\item \textbf{Tratar a data de disponibilidade como parte do dado.} O relatório de posições de fundos refere-se a uma terça-feira e sai na sexta; os preços de insumos do mês saem dias depois do fechamento. Ancorar cada série na data em que ela ficou conhecida, e não na data a que se refere, evitou vazamentos que só foram detectados por auditoria explícita do pipeline.
	\item \textbf{Comparar com o que o usuário já tem.} A curva de futuros é a expectativa gratuita do produtor com acesso à bolsa. Uma ferramenta que não a supera não agrega valor preditivo, ainda que supere baselines estatísticos, e deve dizer isso ao usuário em vez de esconder atrás de um RMSE.
	\item \textbf{Validar por \textit{backtest} rolante em mais de um regime.} Vinte decisões em duas safras de alta não distinguem uma política boa de uma sortuda. Dobrar a amostra com um período de queda mudou a leitura de várias estratégias, inclusive da venda parcelada, e é o mínimo para qualquer afirmação de ganho.
	\item \textbf{Dimensionar para o computador do usuário.} O motor resolve uma decisão em milissegundos e o treino da rede leva pouco mais de um minuto em CPU. Cooperativas não precisam de infraestrutura de nuvem para operar a ferramenta, e o retreino, que neste estudo foi feito uma única vez, pode ser executado a cada safra ou a cada mês em um servidor comum, com a ressalva de que cada retreino exige repetir a verificação com várias sementes.
	\item \textbf{Começar pela rede menor.} A configuração inicial, de 64 unidades, parecia modesta e estava sobreajustada a 3\,940 janelas de treino; a busca mostrou que metade das unidades e o dobro do \textit{dropout} melhoram o erro de seis meses em 19\,\% na média de cinco sementes. Em séries agrícolas, com poucas safras de histórico, a capacidade da rede deve ser dimensionada pela amostra, não pela convenção da literatura, e a busca de hiperparâmetros deve ser verificada com várias sementes, porque uma semente favorável pode exagerar o ganho.
\end{alineas}

Uma última observação diz respeito à origem dos parâmetros econômicos. A tarifa de armazenagem e a quebra técnica foram estimativas do desenvolvimento, sem documento de referência, e a análise de sensibilidade mostrou que a decisão é insensível a elas na faixa plausível. Esse é um resultado útil para quem constrói ferramentas semelhantes: nem todo parâmetro precisa ser levantado com precisão antes de a ferramenta ser útil, e a análise de sensibilidade é a forma de descobrir quais precisam.
